\section{Introduction}\label{sec:introduction}\justifying{}
Hypertension, commonly known as high blood pressure, is a chronic cardiovascular condition in which the force of the blood against the artery walls is consistently too high. Uncontrolled hypertension significantly increases the risk of serious health problems, including heart attack, stroke and kidney disease. The renin-angiotensin-aldosterone system (RAAS) regulates blood pressure through a negative feedback mechanism where reduced renal blood flow indirectly causes the angiotensin-converting enzyme (ACE) transforming angiotensin I into angiotensin II, a potent vasoconstrictor that raises blood pressure. Angiotensin II triggers aldosterone secretion which increases sodium and water retention to maintain vascular volume and resistance. Dysfunction of the RAAS system causes hypertension, which can be managed with ACE inhibitors. These widely prescribed drugs lower blood pressure and minimize cardiovascular risk by blocking the conversion of angiotensin I to angiotensin II. Additionally, ACE inhibitors prevent the breakdown of bradykinin, a vasodilator that further reduces blood pressure and alleviates heart strain. Common examples include enalapril, lisinopril, ramipril, captopril, and benazepril, each playing a vital role in cardiovascular management~\cite{Vertes1992}.

Enalapril is an FDA-approved ACE inhibitor drug for treating hypertension. Enalapril is the esterified prodrug, usually given orally as enalapril maleate, and further hydrolysed in the liver (\textit{in vivo} de-esterification) to enalaprilat, the active metabolite involved in ACE inhibition. Enalaprilat acts as a competitive inhibitor of ACE, preventing angiotensin I from interacting with it. The resulting enzyme-inhibitor complex has a slower dissociation rate, which effectively causes ACE inhibition. 

Pharmacokinetic parameters such as maximum concentration (C\textsubscript{max}), time to peak concentration (T\textsubscript{max}), half-life (T\textsubscript{1/2}), and elimination rate (k\textsubscript{el}) inform dosing strategies that maximize efficacy while minimizing side effects. In healthy adults, enalapril has an oral bioavailability of 60\%, unaffected by food~\cite{Swanson1984}, and reaches peak plasma concentrations 1--2 hours after ingestion~\cite{Dickstein1987, MacFadyen1993, Ulm1983, VanHecken2020}. It is subsequently hydrolyzed in the liver by carboxylesterase 1 (CES1) to its only active metabolite, enalaprilat: by 4 hours post-dose, plasma enalapril levels are negligible while enalaprilat approaches its own maximum concentration, reached 3--4 hours post-biotransformation, with an independent bioavailability of approximately 40\%~\cite{Davies1984, Till1984}. Enalapril has an elimination half-life of 2--3 hours and is excreted in urine and feces, and evening administration delays its T\textsubscript{max}~\cite{Weisser1991}. Hypertension itself does not alter the pharmacokinetics of enalapril or enalaprilat~\cite{Todd1992b}, indicating that the disease being treated is not a source of pharmacokinetic variability.

Instead, this variability arises from patient- and organ-level factors, most notably age~\cite{Hockings1986, Lees1987a, MacDonald1993}, renal impairment~\cite{Fruncillo1987, Kelly1986, Lowenthal1984}, and hepatic impairment~\cite{Baba1990, Ohnishi1989}. The evidence for an age effect is mixed: Hockings et al. 1986~\cite{Hockings1986} and Lees et al. 1987~\cite{Lees1987a} found no age-related differences in pharmacokinetics or ACE inhibition between young and elderly groups, whereas MacDonald et al. 1993~\cite{MacDonald1993} reported a higher enalaprilat area under the curve (AUC) in older individuals, indicating prolonged drug exposure. This inconsistency suggests that age alone is an unreliable proxy for the underlying physiological changes that actually drive variability in enalapril pharmacokinetics, motivating a mechanistic treatment of the organ functions -- renal and hepatic -- that age-related decline, renal disease, and liver disease each affect directly.

Because pharmacokinetic exposure ultimately determines pharmacodynamic effect, the same sources of variability are expected to propagate to enalapril's action on the RAAS. For the ACE inhibitor enalapril, key pharmacodynamic parameters include blood pressure, heart rate, ACE activity, and RAAS hormone concentrations (renin, aldosterone, and angiotensins) ~\cite{Donnelly1987, Lees1987a, Ribeiro1996, Shionoiri1985}. Studies demonstrate that a 6-week daily enalapril regimen significantly reduces systolic and diastolic blood pressure, achieving near-complete ACE inhibition (\(101\%\)--\(103\%\)) ~\cite{Donnelly1990}. ACE inhibition typically peaks at 90\% within 3--5 hours post-dose and drops to 50\% at 24 hours ~\cite{Ribeiro1996}. When comparing age groups, younger and older patients show similar levels of ACE inhibition ~\cite{Lees1987a}. While older patients experience a greater absolute drop in blood pressure, this difference disappears when normalized as a percentage, meaning the larger decrease is simply due to their higher baseline blood pressure~\cite{Lees1987a}.

Since enalapril and enalaprilat are eliminated almost exclusively via the kidneys, renal dysfunction directly alters their pharmacokinetics: progressive renal impairment reduces urinary elimination of enalapril, raising peak serum enalapril levels and increasing its conversion to enalaprilat~\cite{Fruncillo1987}, so that patients with renal failure show higher serum enalaprilat concentrations, longer T\textsubscript{max}, and lower excretion rates for both the parent drug and its metabolite~\cite{Kelly1986, Lowenthal1985}. Because enalapril additionally depends on hepatic CES1-mediated conversion to become active, liver dysfunction represents a second, independent route by which organ function shapes exposure -- though here the evidence is less consistent. Ohnishi et al.~\cite{Ohnishi1989} found that cirrhotic patients had higher enalapril C\textsubscript{max} and AUC alongside lower urinary excretion, indicating impaired hepatic processing, whereas Baba et al. 1990~\cite{Baba1990} reported no significant differences in C\textsubscript{max}, T\textsubscript{max}, or ACE inhibition between cirrhotic and healthy subjects. Renal and hepatic function thus each independently, and potentially jointly, determine enalapril exposure, motivating a model that can separate and combine their effects rather than treating impairment as a single covariate.

Genetic variation in CES1 provides a third, molecularly defined source of variability. The loss-of-function G143E variant severely impairs enalapril metabolism: carriers show an approximately 20--27.5\% reduction in enalaprilat AUC, lower peak concentrations, and a roughly 30\% delay in T\textsubscript{max}~\cite{Her2021, Tarkiainen2015}, with a correspondingly diminished therapeutic effect -- a negligible systolic blood pressure drop (-1.0 mmHg) compared to 14.6 mmHg in wild-type patients~\cite{Her2021}. Stage et al. 2017~\cite{Stage2017} evaluated additional complex variants, including copy-number variations (3-copy and 4-copy) and structural changes (CES1A2, CES1A1c); while several reduced enzymatic activity or prolonged T\textsubscript{1/2}, median enalaprilat AUC across these variants did not differ significantly from wild type. CES1 genotype therefore modulates enalapril pharmacokinetics through the same hepatic conversion step affected by cirrhosis, but via a distinct, genetically defined mechanism that a purely organ-function-based model cannot capture on its own.

With this study, we aim to develop a reproducible, openly available PBPK/PD model of enalapril, built from published data in healthy subjects and patients with hypertension, renal impairment, and hepatic impairment. This allows comparative analysis across health states to understand how hepato-renal function and CES1 activity affect enalapril's distribution, metabolism, and excretion. Because enalapril's effects vary with patient characteristics, organ impairment, and CES1 genetic variation, the model provides a framework for predicting drug behavior across these conditions -- supporting more personalized dosing strategies and filling gaps in current mechanistic understanding.